%=====================================================================
%  Front matter: abstract, contents, lists, abbreviations
%=====================================================================
\chapter*{Abstract}
\addcontentsline{toc}{chapter}{Abstract}

Equity markets are driven simultaneously by structured price and volume records, by
qualitative information carried in news and social media, and by slower macroeconomic
conditions. Each stream is individually incomplete and each becomes informative or
uninformative at different times. Conventional forecasting systems compress this heterogeneity
into a single numerical pipeline, report a statistical accuracy figure, and stop short of the
decision the investor actually has to take. This research addresses that shortfall by
developing, implementing and validating an integrated framework in which multi-source fusion,
language-model-based sentiment extraction, transparent reasoning and decision-level allocation
are treated as one connected engineering problem rather than as four separate ones.

The first objective develops a hierarchical Multi-Source Fusion Network combining
convolutional indicator encoders, transformer text encoders, engineered macroeconomic features
and a cross-modal attention layer that assigns time-varying weight to each source. On ten
years of National Stock Exchange data across ten sectoral indices it attains a normalised mean
absolute error of 0.018, $R^2$ of 0.988 and directional accuracy of 84.7 per cent at the index
level, against 76.3 per cent for a simple averaging ensemble over identical inputs; removing
both textual sources costs 10.9 per cent of explained variance, exceeding the sum of their
individual removals. The fusion rule is then strengthened by a Bayesian formulation in which
the weight on a source is derived from its own recent predictive uncertainty through a softmax
over negative error variances, improving the Sharpe ratio by 22.7 per cent and reducing
maximum drawdown by 18.3 per cent against the strongest static-fusion baseline, while the same
experts fused with fixed weights add nothing --- which locates the gain in the weighting
mechanism rather than in the networks.

The second objective couples a transformer sentiment encoder with fifteen technical indicators
under a strict temporal protocol. Language-model sentiment emerges as the highest-gain
predictor, correlating 0.74 with momentum yet only 0.04 with price, and raises directional
accuracy from a 46.43 per cent baseline mean to 47.72 per cent while producing the highest
prediction confidence of all models examined. The absolute level is reported without
embellishment: at single-name daily horizons all models cluster near the random classifier. The
third objective replaces correlation-based feature selection with a directed causal adjacency
matrix recovered by conditional independence testing, encodes intraday, daily and weekly
horizons through a multi-scale memory encoder, and allocates capital through a credibilistic
optimiser. Over 308 NIFTY50 sessions it recovers a causal graph of density 0.850 with an
identifiable causal hub and shows that strong causal influence can coexist with weak linear
correlation; the sentiment composite carries the expected sign against four independently
published series and separates the identified market states at $p = 1.8 \times 10^{-8}$, yet is
separately shown not to forecast returns.

The fourth objective evaluates the framework as a decision system. Over eighteen years of ten
liquid multi-asset instruments under a walk-forward protocol with transaction costs it attains a
Sharpe ratio of 0.878 against 0.628 for an equal-weighted benchmark and reduces maximum drawdown
from $-36.1$ to $-19.0$ per cent, winning in all four stress episodes, in 18 of 19 calendar
years on drawdown and in all 18 specifications examined, and beating all three benchmarks on a
665-session holdout that postdates every design decision. Ablation at matched average exposure
isolates the source of the gain: timing risk by market state improves both the Sharpe ratio and
the drawdown with no change in mean return, whereas the signal fusion layer contributes nothing
detectable. Cross-market transfer reproduces the defensive behaviour but loses on risk-adjusted
return, bounding applicability to universes with genuine dispersion in asset volatility. Taken
together, the six publications underlying this report demonstrate a coherent progression from
information integration, through language-model-based signal extraction and transparent
reasoning, to validated decision making; the component-level nulls are retained because they
identify precisely which layers of a fusion architecture earn their place and which do not.

\vspace{0.3em}
\noindent\textbf{Keywords:} multi-source data fusion; sentiment analysis; large language
models; explainable artificial intelligence; regime detection; portfolio construction; Indian
equity market.

%---------------------------------------------------------------------
\clearpage
\renewcommand{\contentsname}{Table of Contents}
{\setstretch{1.0}\tableofcontents}

\clearpage
\renewcommand{\listfigurename}{List of Figures}
{\setstretch{1.0}\listoffigures}
\vspace{1.2em}
\renewcommand{\listtablename}{List of Tables}
{\setstretch{1.0}\listoftables}

%---------------------------------------------------------------------
\clearpage
\chapter*{Abbreviations}
\addcontentsline{toc}{chapter}{Abbreviations}

\begin{center}
\footnotesize
\setstretch{1.0}
\begin{tabular}{@{}p{1.9cm}p{4.5cm}p{1.9cm}p{4.9cm}@{}}
ANOVA & Analysis of Variance & MAE & Mean Absolute Error\\
ARIMA & Autoregressive Integrated Moving Average & MDD & Maximum Drawdown\\
AUC & Area Under the Curve & ML & Machine Learning\\
Bi-LSTM & Bidirectional Long Short-Term Memory & MLP & Multi-Layer Perceptron\\
BSE & Bombay Stock Exchange & MOSPI & Ministry of Statistics and Programme Implementation\\
CAGR & Compound Annual Growth Rate & MSFN & Multi-Source Fusion Network\\
CIT & Conditional Independence Test & NLP & Natural Language Processing\\
CNN & Convolutional Neural Network & NSE & National Stock Exchange\\
CPI & Consumer Price Index & OHLCV & Open, High, Low, Close, Volume\\
CVaR & Conditional Value at Risk & QP & Quadratic Program\\
DA & Directional Accuracy & RBI & Reserve Bank of India\\
DCHF & Dynamic-Causal Hybrid Framework & RDM-DEA & Range Directional Measure Data Envelopment
Analysis\\
DL & Deep Learning & RMSE & Root Mean Square Error\\
ETF & Exchange-Traded Fund & RSI & Relative Strength Index\\
GMM & Gaussian Mixture Model & SEBI & Securities and Exchange Board of India\\
GRU & Gated Recurrent Unit & SHAP & Shapley Additive Explanations\\
IC & Information Coefficient & SMA & Simple Moving Average\\
IIP & Index of Industrial Production & SVM & Support Vector Machine\\
LLM & Large Language Model & VADER & Valence Aware Dictionary and sEntiment Reasoner\\
LSTM & Long Short-Term Memory & VIX & Volatility Index\\
MACD & Moving Average Convergence Divergence & WPI & Wholesale Price Index\\
& & XAI & Explainable Artificial Intelligence\\
\end{tabular}
\end{center}
